\documentclass[11pt]{article}
\usepackage[T1]{fontenc}
\usepackage{lmodern,amsmath,amssymb,booktabs,longtable,array}
\usepackage[margin=1in]{geometry}
\usepackage[hidelinks]{hyperref}
\hypersetup{pdftitle={The Antidote: Project Overview},pdfauthor={The Antidote project}}
\setlength{\parindent}{0pt}
\setlength{\parskip}{0.65em}
\setlength{\emergencystretch}{2em}
\title{The Antidote\\[0.4em]\large Project overview}
\author{Models, data, and resilience}
\date{}
\renewcommand{\abstractname}{At a glance}
\begin{document}
\maketitle
\begin{abstract}
The Antidote estimates when a ship will cross an analytical harbor gate in Hampton Roads. The historical replay compares quadratic regression, ordinary XGBoost, and Antidote's XGBoost model with conditional input repairs as AIS reports are delayed, missing, noisy, or biased. When an issue is detected, a human records its error type; those reviewed examples are saved and feed the next model-training cycle. The intended user is a port logistics coordinator deciding whether an arrival-dependent plan needs verification.

This overview brings together the full-March nominal comparison with attack and mitigation results. It explains how reports are disrupted, how the models respond, and how to interpret the findings.
\end{abstract}
\setcounter{tocdepth}{1}
{\setlength{\parskip}{0pt}\tableofcontents}

\section{The problem}
The intended user is a port logistics coordinator deciding whether to keep an arrival-dependent plan or verify it before committing resources. A late or unstable estimate can cause avoidable waiting, premature staffing, or repeated replanning. The prototype supplies an estimate and evidence for that decision; it does not allocate a berth, issue navigation instructions, or control a vessel. The \href{/tracking/}{historical replay} supplies the interactive view of ships, incoming reports, data changes, and forecasts. Here, we explain the inputs and decisions behind those forecasts.

Our outcome is an \emph{analytical harbor-gate crossing}, not docking, cargo delivery, or an official port-arrival timestamp. A qualifying east-to-west crossing intersects longitude $\lambda_g=-76.3^\circ$ between latitudes $36.95^\circ$ and $37.05^\circ$. Consecutive original observations must bracket it within 120 seconds. If the two longitudes are $\lambda_0,\lambda_1$, the reference is interpolated as
\[
 \alpha=\frac{\lambda_g-\lambda_0}{\lambda_1-\lambda_0},\qquad
 T^*=t_0+\alpha(t_1-t_0).
\]
Interpolation assumes approximately linear motion inside the bracket. Sampling, position error, and uncertain destination intent remain limitations of the reference. The original bracket is fixed before introducing faults and never becomes a repair input. Cohort membership requires an observed qualifying crossing, so results cannot establish performance for unresolved approaches or vessels that never arrive.

We assess forecasts using MAE, MSE, 95th-percentile absolute error, prediction coverage, review burden, and condition-level results. Together, these measures show average and tail performance, forecast availability, and intervention needs.

\subsection{The data and evaluation window}
\subsubsection{Source and target}
The current serving bundle contains frozen A, B, and nominal-only C models and a replay of all 31 days of March 2024. The matching nominal comparison reports 5,382 scheduled decisions across 243 qualifying approaches and 219 vessels. Its current fit partition is 1 January--29 June 2023 (27,263 decisions, 1,319 approaches); model selection/early stopping uses 30 June--19 July (3,108 decisions, 140 approaches). The model and source hashes record the data lineage and configuration.

The March source is NOAA's cleaned broadcast-point AIS product, selected for Hampton Roads and filtered to cargo and tanker reports. The bundle records the March source checksum and model/code signatures. Its target is an interpolated analytical gate crossing, not verified docking or official arrival. NOAA's metadata describes collection and processing limits.\footnote{\href{https://www.fisheries.noaa.gov/inport/item/73064}{NOAA, Nationwide Automatic Identification System 2024.}}

Shared preprocessing selects cargo and tanker reports before simulated masking, parses UTC times, treats invalid navigation values as missing, and removes exact duplicate measurements consistently across models. MMSI identifies repeated vessel observations; it is not an independent observation count. Source archives already involve processing and incomplete receiver coverage, so ``original'' means unchanged by our simulator, not error-free ground truth.

\subsubsection{Temporal split and information available at inference}
The temporal separation is the study's central data contract. The current frozen models are fit and selected on 2023 data, and evaluated on the same March 2024 cohort. This is not an 80/10/10 random row split. Model settings are selected within 2023, with complete approaches kept together and the March evaluation cohort kept separate from fitting.
\begin{longtable}{@{}p{0.25\linewidth}p{0.25\linewidth}rr@{}}
\toprule
Partition & Dates & Approaches & Clocks \\
\midrule
\endhead
Current fit & 1 Jan--29 Jun 2023 & 1,319 & 27,263 \\
Current model selection & 30 Jun--19 Jul 2023 & 140 & 3,108 \\
Nominal evaluation & 1--31 March 2024 & 243 & 5,382 \\
\bottomrule
\end{longtable}
The nominal evaluation uses 219 vessels across the 243 approaches. The paired A/B/C decision identities match, and the saved test predictions are finite and available. Recurring vessel identity across years means chronological separation does not establish vessel independence. March has also been examined during development, and the choice of nominal C followed that examination; March is therefore not a fresh confirmation set for that model choice.

For report $r_i$, let $t_i$ be its measurement time and $d_i$ its simulated delivery time. At decision clock $t$, inference receives only
\[
 \mathcal D(t)=\{r_i:t_i\leq t,\ d_i\leq t\},\qquad
 a(t)=t-\max_{r_i\in\mathcal D(t)}t_i.
\]
The second quantity is the age of the latest measurement. Nominal delivery is instantaneous by assumption; it is not a measurement of actual receiver latency. Forecasts are scheduled every five minutes after at least ten minutes of an approach. Features summarize the ten minutes ending at the latest available observation, while age continues to increase at the decision clock. This preserves what an operator could know as time advances.

In the live replay, the reporting-gap control withholds reports for a chosen interval; it does not deliver old messages later. A separate saved stress curve tests delayed delivery. Neither substitutes for measured receiver latency.

\section{The solution}
\subsection{Quadratic regression: a structured comparator}
Model A expands the base features with degree-two polynomial terms and fits ridge regression. With feature expansion $\phi(x)$, it minimizes
\[
 \sum_i w_i\bigl(y_i-\beta^\mathsf T\phi(x_i)\bigr)^2
 +\lambda\lVert\beta\rVert_2^2.
\]
The base representation includes local east/north position, current speed, sine and cosine of course, recent mean speed and position, speed trend, history support, and explicit cargo/tanker/unknown-class flags. Circular course encoding avoids treating $359^\circ$ and $1^\circ$ as nearly a full rotation apart. Polynomial interactions provide a compact nonlinear comparator; regularization limits coefficient growth. This model is inspectable, but it can extrapolate poorly during waiting, turning, or unfamiliar approach patterns.

\subsection{Ordinary XGBoost: a nonlinear comparator}
Model B uses gradient-boosted regression trees on the same base feature contract and nominal training inputs. Predictions combine sequential tree contributions, $f(x)=\sum_{m=1}^{M}h_m(x)$. XGBoost supplies the nonlinear learner; it is the AI capability under evaluation, rather than merely a tool used to author the website.\footnote{\href{https://arxiv.org/abs/1603.02754}{Chen and Guestrin, XGBoost: A Scalable Tree Boosting System, 2016.}} Missing numeric features pass through an imputer fitted on training data. Unknown vessel class has its own flags and is not silently turned into cargo.

Both A and B predict remaining duration from the latest observation. Their absolute arrival prediction is $\widehat T_{A,B}=t_{\mathrm{obs}}+f_{A,B}(x)$. Subtracting the current decision clock yields the remaining time now. An older observation can retain the same absolute forecast while its remaining duration shrinks; changing a display clock does not refresh the underlying evidence.

\subsection{Antidote: degradation-aware learning and conditional repairs}
Model C is also XGBoost. It adds robust motion summaries and observable freshness/quality signals: speed quantiles and variability, turn and course dispersion, a position-history velocity fit, report age, gap length, missing-field fractions, occupied minute bins, and disagreement between reported and position-derived motion. It predicts from the decision clock, $\widehat T_C=t+g(z)$. The model served in the current comparison is the nominal-only C artifact: it was fitted on original reports, without injected deterioration. The separate augmented C remains an experimental artifact and is not served. Conditional repairs are added by the dashboard at inference time.

The trained guard limits report influence by forming one representative per UTC minute using medians and a circular course summary. It screens invalid coordinates and motion exceeding 60 knots across short gaps of at most 300 seconds. Minute binning limits the contribution of a prolific feed to history summaries; it does not make a biased majority within a minute trustworthy. Exact deduplication is shared with the ordinary comparators, so it cannot be credited exclusively to Antidote.

The current dashboard adds conditional repair strategies to this frozen C model. These are separate from its learned guard. The code keeps raw A/B inputs, original quality signals, support, labels, and clocks intact while changing only permitted derived C inputs. Training uses squared-error regression; depth and learning rate are selected using approach-balanced development MAE. The configured search spans depths 2, 4, and 6, learning rates 0.03 and 0.1, up to 500 boosting rounds, and early stopping after 30 rounds. Choosing a model by MAE does not ensure its MSE or worst errors improve.

\subsection{Degradation scenarios}
The replay studies accidental deterioration at \emph{inference time}: a frozen model version receives worse reports. No live AIS broadcasts are altered, and a reviewed label does not silently change the model serving that replay. Instead, reviewed cases are retained for the next versioned training cycle. Although controls are called attacks in code, injected changes stand in for outages, unavailable telemetry, miscalibration, and navigation faults; they are not evidence of malicious activity in NOAA records.

\begin{longtable}{@{}p{0.22\linewidth}p{0.34\linewidth}p{0.36\linewidth}@{}}
\toprule
Fault & Controlled change & Mechanism affecting the forecast \\
\midrule
\endhead
Reporting gap & Withhold all reports in a fixed interval beginning ten minutes into an approach. & Old position and motion remain the latest evidence; data age increases while the ship continues moving. \\
Missing speed & Mask selected SOG values. & Current speed, recent summaries, and missingness evidence change. Zero speed remains a legitimate observation. \\
Missing course & Mask selected COG values. & Direction and turning features lose support; numerical imputation may follow a different tree path. \\
Missing vessel class & Mask selected type reports. & The model loses cargo/tanker context; unknown remains a distinct state. \\
Speed drift & Add a slowly varying bounded speed bias. & A persistent underestimate can resemble slower progress and shift predicted duration. \\
Location drift & Accumulate a bounded coordinate displacement. & Distance to the gate and position-derived motion become biased together. \\
Sensor deterioration & Add zero-mean Gaussian noise to SOG and COG. & Instantaneous errors propagate into direction, speed variability, and nonlinear tree decisions. \\
Position jump & Displace selected reports along a seeded direction. & Apparent progress or implausible motion changes position and quality features. \\
\bottomrule
\end{longtable}

For illustration, the navigation perturbations can be written as
\[
 \widetilde v_i=\max\{0,v_i+b_v(t_i)+\epsilon_i\},\qquad
 \widetilde{\mathbf p}_i=\mathbf p_i+\mathbf b_p(t_i)+J_i\mathbf u_i.
\]
Here $b_v$ is persistent speed bias, $\epsilon_i$ is random speed noise, $\mathbf b_p$ is coordinate drift, and $J_i\mathbf u_i$ is a position jump. The simulator also wraps noisy course into $[0,360)^\circ$. Drift and noise differ: averaging can reduce zero-mean jitter under appropriate assumptions, while a persistent bias generally survives averaging.

The saved March graph comparisons use 25\% of each control's configured maximum: a 15-minute reporting gap and 25\% for each other control. The replay sliders affect live replay inputs independently; the graph selector uses the saved fixed-strength scenarios described in the results section.

Faults compose in a fixed order: gap, speed drift, location drift, noise, position jumps, and missing fields. Missing-field and position-jump selectors share a deterministic report hash and seed. Equal-severity selections therefore overlap subject to field eligibility; they are not independent dropout masks. Modification logs are constructed separately. Audit labels, original altered values, and future crossing outcomes are exposed to the audience only after inference.

\subsection{Repairs, warnings, and deferral}
\subsubsection{Reconstructing missing motion and noisy sensors}
When positions are sufficiently consistent, a short local track can estimate motion. For local east/north position $\mathbf p_i$ and times centred around $\bar t$, the least-squares velocity is
\[
 \widehat{\mathbf v}=
 \frac{\sum_i(t_i-\bar t)(\mathbf p_i-\bar{\mathbf p})}
      {\sum_i(t_i-\bar t)^2},\qquad
 \widehat s=\frac{\lVert\widehat{\mathbf v}\rVert}{0.514444},\qquad
 \widehat\theta=\operatorname{atan2}(\widehat v_E,\widehat v_N).
\]
Positions are in metres and time is in seconds; $\widehat s$ is therefore in knots. Course is converted to degrees and wrapped around the circle. A stable trajectory can support estimates for missing speed and course independently. Both use original received positions; one reconstructed sensor does not become evidence for the other. Each historical correction uses only information available when that particular report arrived, preventing later reports from backfilling an earlier estimate.

The Gaussian-noise strategy checks a five-minute track with at least three distinct reports, at least 60 seconds of span, no gap above 180 seconds, and fit residual at most 50 metres. Additional speed, turn, and plausibility checks reject unsuitable references. It replaces a supported reading when its speed differs by more than 2 knots or its moving-course difference exceeds $15^\circ$. This is \emph{guarded replacement}, not a fitted blend of the reported and estimated readings. A smoother track is still only a conditional reference, not independent ground truth.

\subsubsection{Detecting persistent bias}
Speed drift is checked against position-supported speed. For disagreement $r_k=v_k-\widehat s_k$, the detector forms a time-aware exponentially weighted estimate,
\[
 e_k=(1-\rho_k)e_{k-1}+\rho_k r_k,\qquad
 \rho_k=1-\exp(-\Delta t_k/\tau),\qquad \tau=2\text{ minutes}.
\]
The provisional rule requires same-direction disagreement across at least three reports and two minutes, with $|e_k|>0.5$ knots. Reference acceleration, track fit, and turning checks can prevent reconstruction. These conditions reduce impulsive replacements but also reject legitimate manoeuvres; the screen finds limited repair availability for speed drift.

Location drift compares observed displacement with motion implied by received SOG/COG. Persistent displacement disagreement can justify verification, but it cannot identify the true coordinates by itself. The dashboard therefore warns rather than correcting location. A separate experimental mode projects from explicit, time-available position anchors for at most five minutes and requires trusted navigation. Those anchors are simulated from original AIS in the experiment, so they are an oracle-reference study, not an independent operational feed. Coherent corruption of both position and motion may evade these consistency checks.

\subsubsection{Fallbacks, metadata, and combined faults}
The reporting-gap strategy detects silence from observed cadence and does not fabricate reports. Position-jump detection supplies navigation-consistency warnings and does not repair coordinates. Missing class can be carried forward only from consistent recent received metadata for the same MMSI; conflicts remain unknown and require review.

Position faults pause position-based navigation repairs. Noise plus speed drift pauses their overlapping sensor corrections; class recovery can continue independently. The current fleet policy selects permitted reference assumptions from the simulated fault families enabled in the controls. It is not blind automatic fault isolation. Pausing avoids asserting an unsupported correction, but it is not evidence that the forecast has recovered. A production system would need independently justified reference trust and an evaluated fallback policy.

\section{The results}
\paragraph{Full-March nominal baseline.}
The serving bundle scores A, B, and nominal-only C on the same 5,382 scheduled decisions across 243 approaches and 219 vessels. These are original, unmodified March reports; no injected faults or additional dashboard repairs are applied in this comparison. MAE is averaged within each approach, then across approaches. MdAE and P95 are pooled across individual forecast errors.
\begin{table}[!htbp]
\centering\small
\begin{tabular}{@{}lrrrr@{}}
\toprule
Model & Approach-weighted MAE (min) & Forecast MdAE (min) & Forecast P95 AE (min) & Coverage \\
\midrule
A: quadratic ridge & 7.019 & 3.459 & 36.725 & 100\% \\
B: ordinary XGBoost & 6.643 & 3.164 & 37.361 & 100\% \\
C: nominal-only XGBoost & 6.610 & 3.097 & 37.097 & 100\% \\
\bottomrule
\end{tabular}
\end{table}
C's MAE is 0.033 minutes (about 2 seconds) below B's. The paired vessel-bootstrap 95\% interval for C minus B is $[-0.125,+0.050]$ minutes and spans zero; this does not establish a reliable population-level advantage. A has the lowest P95 error. These are point predictions before any interval-based deferral. March was held out from model fitting and tuning, but the choice of nominal C followed inspection of March results, so March is not an untouched confirmation set for that choice.

\subsection{Staleness}
At each decision clock, the model receives reports only after their simulated delivery time, and report age continues increasing until a newer measurement arrives. The live replay's gap control withholds reports, while a separate saved stress curve recomputes forecasts under delayed delivery; these are distinct scenarios. Nominal delivery is assumed instantaneous, so neither is a measurement of real receiver latency. The baseline table above gives the full-March reference; stress graphs show how forecasts move as delayed evidence accumulates.

\subsection{Low-quality data and source volume}
The replay separately exposes missing fields, noisy or drifting sensors, and position jumps, then allows these conditions to be combined. Fault definitions and affected fields are shown above. The current scenario bundle uses a fixed 25\% strength and six sampled approaches; it does not provide a source-volume sweep that increases the share or repetition of low-quality reports. We therefore do not claim to have measured when a higher-volume unreliable source outweighs a lower-volume reference.

\subsection{Over-time behavior}
The ordered replay advances through received reports and forecast updates. Saved comparisons cover all 256 on/off subsets of the eight controls beginning at the March 15 cutover, using six of the 243 March approaches selected by time and duration without looking at prediction errors. The graphs do not recalculate when live sliders move. Their saved predictions have source/model/code signatures and 108 scalar forecast checks against the serving path, but the six-approach sample has no fleet-level uncertainty interval.

\subsection{Countermeasure effectiveness and limits}
The live view shows Antidote's conditional repair, the affected report fields, and whether the available evidence supports correction or requires review. The full-March nominal comparison above does not apply additional dashboard repairs. A same-case comparison of C with repairs disabled and enabled remains to be completed across nominal, stale-only, low-quality-only, and combined conditions. That evaluation must also quantify normal-data effects, coverage, large errors, and performance on a later or otherwise unseen condition. Until then, the current graphs explain scenarios and repair behavior; they do not establish a fleet-wide mitigation gain.

Across the product, the reference remains an observed analytical crossing, not verified docking. The simulator represents assumed data faults and is not evidence of malicious activity in NOAA records.

For approach $j$ with $n_j$ available forecasts and signed error $e_{jk}=\widehat T_{jk}-T_j^*$ in minutes, the primary accuracy metrics are
\[
 \operatorname{MAE}=\frac1J\sum_{j=1}^J\frac1{n_j}\sum_k|e_{jk}|,
 \qquad
 \operatorname{MSE}=\frac1J\sum_{j=1}^J\frac1{n_j}\sum_k e_{jk}^2.
\]
Coverage uses all scheduled cases as its denominator, including unavailable forecasts. Large-error percentiles pool clocks. Paired uncertainty should resample vessels, retain repeated approaches within a vessel, and be reported alongside condition-level outcomes. The current full-March nominal comparison uses a paired vessel bootstrap; the six-approach attack graphs do not supply uncertainty intervals for a fleet-level claim.

\section{Where humans are involved}
When Antidote detects a data-quality issue or an estimate that warrants attention, the reviewer sees the received timestamp, affected fields, diagnostic evidence, and available forecasts. The reviewer identifies the error type when the evidence supports one (for example, a stale report, missing field, sensor noise, speed drift, location drift, or position jump); uncertain cases remain unresolved rather than being forced into a category. The reviewed label and its evidence are saved as a feedback record for later model training. The reviewer does not overwrite the original AIS observation or invent an arrival outcome.

For a reviewed event $i$, let $x_i$ be the evidence available at the decision time and let $z_i$ be its human-adjudicated error type. The feedback store retains the event context and model version alongside $(x_i,z_i)$. A later training cycle adds these reviewed examples to the versioned training set, so a candidate can learn which observable patterns correspond to recurring error types:
\[
 \widehat\phi_{v+1}=\arg\min_{\phi}\left[-\sum_{i\in\mathcal F_v}\log p_{\phi}(z_i\mid x_i)\right].
\]
Here $\mathcal F_v$ contains only reviewed feedback available to training version $v+1$. If a verified crossing outcome $y_i$ is also available, its example can contribute to the ETA objective, $\ell_{\mathrm{ETA}}(f_{\theta}(x_i),y_i)$; an error-type label by itself is not treated as ETA ground truth. This creates a repeatable improvement loop: detect, review, save, retrain, and validate. New candidates are evaluated as versioned models before they replace the frozen replay model, so improvement is the goal of each cycle rather than a guaranteed outcome of every retrain.

An operator should inspect the received timestamp, altered-field evidence, active warning, and available estimates before accepting an unreliable arrival-dependent plan. Simulator audit highlighting explains the experiment and is not used as a human error label or a deployable fault detector. A position inconsistency calls for checking an independent navigation or port source. Antidote exposes conditional estimates without overwriting the displayed received measurements; correction of a source report belongs to the accountable data provider or a separately verified input workflow.

The latest status policy separates an unavailable optional noise/drift reconstruction from an actionable review request. Turning or accelerating can invalidate a simple track fit without proving a fault. Such a case shows ``Estimate unavailable'' while preserving the input evidence. Detected inconsistencies, blocked corrections, missing required support, and other existing review warnings remain. This follows the general principle that alerts should prompt a useful human action.\footnote{\href{https://sre.google/sre-book/monitoring-distributed-systems/}{Google SRE, Monitoring Distributed Systems.}}

Human-reviewed error types and their supporting context are saved for retraining. These labels provide feedback about observed failure modes, but do not by themselves measure alert precision or operator workload. No representative operator study has been completed; a walkthrough should separately record review, verification, correction, and escalation time, plus missed unreliable forecasts and unnecessary interruptions.

\section{Data governance and limitations}
The reviewed NOAA product is publicly accessible; its metadata lists coastal and ocean planning as the use constraint. Public access alone is not blanket authorization for redistribution or operational navigation. Retain attribution, preserve the selected product's metadata, and verify applicable terms before publishing raw subsets or acquiring a live feed. The serving bundle includes model artifacts, a March replay cache, and hashes; source training CSVs are not included in the tracked bundle. The 2023 acquisition record still needs source URLs, retrieval dates, geographic filters, completeness notes, and counts tied to the exact training view.

MMSIs and precise historical tracks identify vessels and can expose operational patterns. We use identity only for grouping and display, without attempting to identify crew, enrich operator profiles, or join unrelated personal information. A live pilot would need explicit access, retention, logging, and sharing rules. Geography, receiving coverage, cargo/tanker selection, repeat vessels, and exclusion of unresolved approaches all constrain representation. Accuracy should be audited by vessel class, forecast horizon, data support, and previously seen vessels before broader claims.

The feedback records retain only the information needed to interpret and reproduce a reviewed event: the relevant input snapshot, assigned error type, model version, review state, and a verified crossing outcome when available. Simulator-generated fault flags remain experiment metadata and do not become human labels. In a proposed pilot, the data steward owns provenance, access, retention, and permitted use; the model lead owns versioned retraining, validation, drift monitoring, and rollback; an operations reviewer owns adjudication and escalation. The submitting team remains accountable for claims and source review. AI assistance was used for implementation and drafting. A language model is not part of the ETA inference path. This overview describes a research prototype, not a certified navigation aid, deployment approval, or a formal human-readiness rating.

\section{Next steps}
The present decision is \textbf{iterate}. The product works as a historical replay, and the nominal March comparison has a saved results and verification record. Current attack evidence is sampled, fixed-strength, and insufficient to establish full-cohort mitigation value or operator readiness.
\begin{enumerate}
\item \textbf{Model lead:} finish the all-eligible-2023 training protocol using time-ordered model selection within 2023. Freeze model and repair choices before the final March scoring pass; disclose the existing March exposure and treat a later period as confirmation. Using the existing offline CPU workflow and matching authorized source files, evaluate every March clock, multiple fault seeds, isolated and combined conditions, coverage, MSE, tail errors, and the same-C repair ablation. Repeat source-volume controls and investigate any large-error cases using development data only; no live AIS integration is needed for this study.
\item \textbf{Data steward:} complete the 2023 and March acquisition manifests, including source URLs, dates, geographic filters, checksums, rights, and retained approach counts. Verify redistribution terms and obtain an independent position reference before making coordinate-recovery claims.
\item \textbf{Demo and operations reviewers:} guide judges through clean, stale, low-quality, combined, repaired, and failed examples. Then run an explicitly exploratory oversight walkthrough that records human effort separately from model latency. Follow the institution's required approval process before any participant study; a peer walkthrough alone cannot establish operator readiness.
\end{enumerate}
The checkpoint is a versioned report showing nominal and degraded accuracy, MAE/MSE and large-error behavior, output coverage, alert misses, review burden, and resource use. Revisit a controlled pilot only if the reach goals hold without unacceptable tail-error or workload tradeoffs, a port partner approves data access, and accountable roles and non-actuating boundaries are in place. Stop or narrow the scope if those conditions fail.

For a ten-minute judge walkthrough, spend one minute on the port-planning decision and crossing reference, three minutes on the clean replay and nominal comparison, three minutes showing a reporting gap and a sensor-quality fault, two minutes on Antidote's repair and a case where it cannot safely correct, then one minute on the measured limits and next evidence needed.


\section{Reproducibility}
The current nominal comparison, model artifacts, March replay cache, hashes, split counts, and inference signatures are saved under \texttt{Data\_Science/artifacts/data\_eval/optimized/} and \texttt{Data\_Science/experiments/nominal\_model\_comparison/}. Its results are reported above as the current full-March nominal comparison. The attack graphs are saved in \texttt{Frontend/tracking/data/scenarios.json}; their metadata records the six selected approaches, fixed strength, source/model/code hashes, and scalar parity checks. Full source CSVs are not included in the tracked bundle, so exact retraining or graph regeneration requires the matching authorized source files.

The serving bundle contains the model artifacts needed for replay; the website does not train on startup. A fresh evaluation run still needs compatible Python dependencies and the matching source inputs. Full setup and preparation commands are in the repository's \href{https://github.com/JeremyBarry03/2026_WM_Hackathon/blob/main/Frontend/README.md}{frontend guide} and \href{https://github.com/JeremyBarry03/2026_WM_Hackathon/blob/main/Data_Science/Optimized_Model/README.md}{model guide}. From \texttt{Data\_Science}, the normal service command is:
\begin{verbatim}
python -m Optimized_Model serve --port 8765
\end{verbatim}
The saved March comparison bundle can be checked with \texttt{python experiments/static\_graphs/export.py --verify-only}. Its source/model/code signatures and scalar forecast checks support consistency with the serving path, not operational correctness. A recorded fallback presentation is a separate submission artifact and is not included here.

\begin{enumerate}
\item \href{https://www.fisheries.noaa.gov/inport/item/73064}{NOAA Office for Coastal Management. Nationwide Automatic Identification System 2024.} Data lineage, fields, constraints, and receiver-coverage limits.
\item \href{https://www.navcen.uscg.gov/how-ais-works}{U.S. Coast Guard Navigation Center. How AIS Works.} Reported position, motion, and vessel-information context.
\item \href{https://arxiv.org/abs/1603.02754}{Chen, T., and Guestrin, C. (2016). XGBoost: A Scalable Tree Boosting System.} Nonlinear learner used in B and C.
\item \href{https://sre.google/sre-book/monitoring-distributed-systems/}{Google SRE. Monitoring Distributed Systems.} Actionable alerting and excessive-interruption risks.
\end{enumerate}
\end{document}
